\documentclass[11pt]{article}
\usepackage[margin=25mm]{geometry}
\usepackage{fontspec}
\setmainfont{Latin Modern Roman}
\usepackage{amsmath,amssymb,mathtools}
\usepackage{unicode-math}
\setmathfont{Latin Modern Math}
\usepackage{xurl,hyperref}
\hypersetup{colorlinks=true,urlcolor=blue}
\setcounter{secnumdepth}{0}
\setlength{\parindent}{0pt}
\setlength{\parskip}{5pt}
\setlength{\emergencystretch}{3em}
\newcommand{\sref}[2]{\hyperlink{source:#1:#2}{[S#2]}}
\newcommand{\cataloguescope}[1]{\paragraph{Recorded target.}#1}
\newcommand{\R}{\mathbb R}
\newcommand{\C}{\mathbb C}
\newcommand{\N}{\mathbb N}
\newcommand{\Z}{\mathbb Z}
\newcommand{\Q}{\mathbb Q}
\newcommand{\F}{\mathbb F}
\newcommand{\eps}{\varepsilon}
\begin{document}
\section{Area–Angular-Momentum Inequalities with Nonpositive Cosmological Constant}\label{statement}

\subsection{Definitions and mathematical statement}

\paragraph{Shared notation.}$\mathbb N=\{1,2,\ldots\}$; the standard number fields have their usual meanings. All additional conventions are stated locally.


\paragraph{Definitions and conventions.}
For an axially symmetric horizon or surface $S$, $A$ is its area and $J_S$ its quasilocal angular momentum in geometrized units. The relevant surface class in the report includes stable minimal surfaces and stable marginally outer trapped surfaces; stability and matter energy conditions constrain admissible geometries. An area--charge inequality compares $A$ with the electromagnetic charges. Equality rigidity identifies all data saturating the sharp bound, rather than merely constructing one equality example.

\paragraph{Target retained from the canonical source catalogue.}
\paragraph{Catalogue wording retained for traceability.}
Extend the area–angular-momentum inequality $A\geq8\pi|J_S|$ to cosmological constant $\Lambda\leq0$, determine the role of extremal anti–de Sitter geometries in the sharp estimate, and prove the equality rigidity statement for the area–charge inequality.

\paragraph{Source scope and qualification.}
The source $A\geq8\pi|J_S|$ is the reference inequality. The question asks for the appropriate $\Lambda\leq0$ extension and extremal AdS role, rather than presupposing that its unchanged right-hand side is sharp in every case. The charge rigidity task is also retained.
\subsection{Short English statement}
Find sharp horizon area bounds with zero or negative cosmological constant and classify equality in the related charge bound.
\subsection{Research attempt}
\paragraph{Model, date and investigated scope.}
OpenAI GPT-6 Astra (configured model \texttt{gpt-6-astra}), ultra reasoning; generated 11 October 2026. Derived the extremal charged-AdS area relation as a normalization benchmark; no general rotating inequality or equality rigidity is established.

\paragraph{Sources and literature check.}
Read the supplied definition, catalogue wording, source qualifications and reference entries for this question. The historical source/context and inherited status are retained; no fresh literature or novelty certification is claimed. No external theorem is used without proof in this attempt; explicit model assumptions are distinguished from facts about the source model.

\paragraph{An extremal charge benchmark.} Test the radial lapse \(f(r)=1-2M/r+Q^2/r^2+r^2/L^2\), \(L>0\), of the spherical charged AdS family. A double horizon at \(r=r_0>0\) requires \(f(r_0)=f'(r_0)=0\). The derivative equation gives \(M=Q^2/r_0-r_0^3/L^2\); substituting it into the first equation yields
\[
 Q^2=r_0^2+3r_0^4/L^2,\qquad M=r_0+2r_0^3/L^2.
\]
For its round area \(A=4\pi r_0^2\), this reads \(Q^2=A/(4\pi)+3A^2/(16\pi^2L^2)\). The extra term is positive, and vanishes in the flat-cosmological limit \(L\to\infty\). Thus the unchanged flat charge-area equality is not the appropriate extremal reference at finite AdS radius.
\paragraph{Gap and audit.} This is an algebraic calculation within a known exact-family ansatz, not a proof that a general stable surface obeys its corresponding inequality. Nor does one equality example imply equality rigidity: that requires every saturating geometry and matter field to coincide with the proposed extremal data. The area--angular-momentum question additionally involves rotation, absent here since this family has \(J=0\). The first missing step is the stability inequality and its sharp minimization with the cosmological term in the actual surface class. Both the rotating extension and charge-rigidity task remain; the benchmark is not claimed as new progress.

\paragraph{Verification and adversarial audit.}
Checked both double-root equations, the flat limit and the distinction between an extremal example and rigidity of all equality cases. This is a GPT-6 Astra Ultra self-audit, not independent human verification. Prompt A drafting followed by prompt B scope/dependency/boundary self-audit. Checked both double-root equations, the flat limit and the distinction between an extremal example and rigidity of all equality cases. The remaining gap is: Obtain and sharply minimize the stable-surface variational inequality with negative cosmological term, including rigidity conditions. Unresolved status and zero original-target completion are retained.

Coordinator checked the two double-root equations and conversion to the extremal area-charge relation. No general surface inequality or rigidity theorem is inferred. This additional check was model-assisted, not independent human verification.

\paragraph{Tools, process and contributions.}
Python 3 and local file-reading/writing tools for manuscript serialization and structural checks; mathematical calculations were symbolic reasoning, with no numerical experiment or theorem prover. Contribution prompt A was followed by prompt B. GPT-6 Astra with ultra reasoning drafted the research and performed its stated self-audit; any additional coordinator contribution is identified in the verification and audit. A GPT-6-based Codex coordinator assembled the manuscript and checked repository integration; its exact variant and reasoning configuration were not exposed. No human mathematical contribution, independent human proof check or formal proof certification is supplied. Complete derivations and the consulted source references are included here; private reasoning transcripts are not published.

\paragraph{Final status and first remaining gap.}
\textbf{UNRESOLVED}. Obtain and sharply minimize the stable-surface variational inequality with negative cosmological term, including rigidity conditions.

\section{Completion Estimate}
\noindent\textbf{COMPLETION: 0\%}
The proved calculation is an elementary auxiliary result or a test of a restricted strategy; no new portion of the full source target is established.

\subsection{Sources}
\begin{itemize}
\item[\hypertarget{source:statement:1}{S1}] ULAM.AI and the UnsolvedMath Contributors, UnsolvedMath, catalogue entry. Catalogue curation under CC BY 4.0; definitions and English exposition added locally.
\url{https://huggingface.co/datasets/ulamai/UnsolvedMath}.
\item[\hypertarget{source:statement:2}{S2}] Mathematical General Relativity (2018), Oberwolfach Reports, DOI 10.4171/owr/2018/36, pp. 2157--2251. Primary report accessed 2026-10-11; relevant PDF pages 78.
\url{https://doi.org/10.4171/owr/2018/36}.
\item[\hypertarget{source:statement:3}{S3}] Gabach Clement--Jaramillo review (catalogue-listed context; not independently consulted in this preparation).
\url{https://doi.org/10.1007/s41114-018-0014-7}.
\end{itemize}
\label{end:statement}
\end{document}
